\subsection{The case of absolute values}

THEOREM 2 (Approximation Theorem for Absolute Values). Let \(f_{i}\) ( \(1 \leqslant i \leqslant n\) ) be absolute values on the same field \(\mathbf{K}\) which are not improper and no two of which are equivalent. Let \(a_{i}\) ( \(1 \leqslant i \leqslant n\) ) be elements of \(\mathbf{K}\) and \(\varepsilon\) a real number \(>0\) . Then there exists \(x \in \mathbf{K}\) such that \(f_{i}(x - a_{i}) \leqslant \varepsilon\) for all \(i\) .

Let \(K_{i}\) denote the field K with the topology defined by \(f_{i}\) . The result to be proved is equivalent to the following: in the product \(P = K_{1} \times \cdots \times K_{n}\) , the closure \(\overline{D}\) of the diagonal D is equal to P. This is obvious for n = 1. Suppose that this point has been established in the case of k absolute values for k < n.

We show first that there exists, for \(2 \leqslant h \leqslant n\) , an element \(x_{h}\) of K such that \(f_{1}(x_{h}) < 1\) , \(f_{2}(x_{h}) > 1\) and \(f_{i}(x_{h}) \neq 1\) for \(3 \leqslant i \leqslant h\) . We argue by induction on h. If h = 2, this follows from the fact that \(f_{1}\) and \(f_{2}\) are not equivalent. We therefore suppose that the existence of \(x_{h-1}\) has been shown and prove that of \(x_{h}\) . If \(f_{h}(x_{h-1}) \neq 1\) , we may take \(x_{h} = x_{h-1}\) ; if \(f_{h}(x_{h-1}) = 1\) , we choose \(z \in K^{*}\) such that \(f_{h}(z) \neq 1\) and \(x_{h} = z(x_{h-1})^{s}\) solves the problem for s sufficiently large. We have thus proved the existence of \(x_{h}\) .

As the integer \(q\) tends to infinity, \(f_{1}(x_{n}^{q})\) tends to 0, \(f_{2}(x_{n}^{q})\) tends to \(+\infty\) and \(f_{i}(x_{n}^{q})\) tends to 0 or \(+\infty\) for \(i \geqslant 3\) . Writing \(y_{q} = x_{n}^{q}(1 + x_{n}^{q})^{-1}\) ,

\[
1 - y _ {q} = (1 + x _ {n} ^ {q}) ^ {- 1};
\]

hence the sequence \((y_q)\) tends to 0 in \(\mathbf{K}_1\) , to 1 in \(\mathbf{K}_2\) and to 0 or 1 in \(\mathbf{K}_i\) for \(i \geqslant 3\) . By changing the numbering of the \(\mathbf{K}_i\) , it may therefore be assumed that there exists an integer \(r\) ( \(1 \leqslant r < n\) ) such that \(\overline{\mathbf{D}}\) contains the point \((e_1, \ldots, e_n)\) where \(e_i = 1\) for \(1 \leqslant i \leqslant r\) and \(e_i = 0\) for \(r + 1 \leqslant i \leqslant n\) . Now, \(\overline{\mathbf{D}}\) is a vector sub-K-space of P. Hence \(\overline{\mathbf{D}}\) contains the diagonals \(\mathbf{D}'\) and \(\mathbf{D}''\) of

\[
\mathbf {P} ^ {\prime} = \mathbf {K} _ {1} \times \dots \times \mathbf {K} _ {r}
\]

and \(P'' = K_{r+1} \times \cdots \times K_n\) . By the induction hypothesis, \(P' = \overline{D}'\) and \(P'' = \overline{D}''\) . Hence \(\overline{D} = P\) .

